
\section{Concept Proposals}
\begin{comment}
    \subsection{Orbital Mechanics}
For a preliminary analysis, we consider basic orbital mechanics to derive minimum fuel requirements for placing a designated satellite into orbit. Orbital transfers and inclination changes require a total change in velocity, denoted as $\Delta V_{tot}$. After \textit{n} burns, the cumulative $\Delta V_{tot}$ is given by:
\begin{equation}
    \Delta V_{tot}= \sum_{i=1}^n |\Delta V_i|
\end{equation}
where each $|\Delta V_i|$ represents the magnitude of the velocity change for each burn \cite{wiesel2003spaceflight}.To transfer from one orbit to another along the same inclination, the required change in velocity, $\Delta V$, is:
\begin{equation}
    \Delta V = V_{f} - V_{i} = \sqrt{\frac{2\mu}{r_2} - \frac{\mu}{a}} - \sqrt{\frac{2\mu}{r_1} - \frac{\mu}{a}}
\end{equation}
where $V_i$ and $V_f$ are the initial and final velocities, $\mu$ is the standard gravitational parameter, $r_1$ and $r_2$ are the radii of the initial and target orbits, and $a$ is the semi-major axis of the transfer orbit \cite{battin1999introduction}. To change the inclination of an orbit, the required change in velocity is:
\begin{equation}
    \Delta V = 2V\sin\left(\frac{\theta}{2}\right)
\end{equation}
where $V$ is the orbital velocity and $\theta$ is the change in inclination angle \cite{curtis2020orbital}. To determine the mass of propellant required for the maneuvers, we apply the Tsiolkovsky rocket equation:
\begin{equation}
    \Delta V = u_e \ln\left(\frac{m_1}{m_2}\right) = I_{sp} g_0 \ln\left(\frac{m_1}{m_2}\right)
\end{equation}
where $u_e$ is the effective exhaust velocity, $I_{sp}$ is the specific impulse, $g_0$ is the standard gravitational acceleration, $m_1$ is the initial mass (including propellant), and $m_2$ is the final mass. Rearranging, the mass of fuel required can be expressed as:
\begin{equation}
    m_{\text{fuel}} = m_1\left(1 - e^{-\Delta V/u_e}\right) = m_1\left(1 - e^{-\Delta V/I_{sg}g_0}\right)
\end{equation}
Note that these calculations assume an instantaneous burn. For preliminary propellant estimation, this assumption provides a reasonable trade-off \cite{wiesel2003spaceflight}.
\end{comment}
\subsection{Concept One: ORACLE}
\begin{center}
        \textit{A look into the past, a glimpse into the future.}
\end{center}

\subsubsection{Abstract}
The first concept follows a low-risk approach, closely resembling the original Hubble Space Telescope design. This is, the satellite shall orbit in Low Earth Orbit (LEO) and equipped with an onboard telescope to capture images of the universe . However, this concept integrates modern instruments to enhance the satellite’s mission capabilities [N1]\footnote{These notations (N1) represents user need reflected.}. The key improvement in ORACLE over the original Hubble is the use of new technology and materials, surpassing the outdated systems available before 1990 (original Hubble launch year) - allowing ORACLE to aligns more closely with the capabilities of the next-generation JWST (greater performance in resolution, better operation efficiency, and lower costs), as well as using more sustainable materials and manufacturing practices [N5]. Additionally, by reflecting the Hubble's past history, literature study, and usage of high Technology Readiness Level (TRL) components from the predecessor, it promises the greatest mission redundancy and theoretically best chance to success. The goal of this concept is to implement these technological advancements within the original orbital trajectory, creating a highly redundant satellite that supports Infrared (IR), Visual (V), and Ultraviolet (UV) imaging tasks - maximizing confidence and minimizing risk with greatest available efficiency [N2].


\subsubsection{Orbit Type} 
 LEO orbit offers the most economical, service accessible, and simple way to engineer an astronomical telescope [N7]. It offers the significant improvement from a ground observatory, considering weather, atmosphere, and importantly UV \& IR attenuation is no longer a factor. With existing communication infrastructure surrounding the Hubble today, ground station regarding the Hubble telescope does not vary with constant orbit trajectory [N4]. The complexity of the satellite's construction/material requirement will also be significantly lowered compared to deep space travel and Lagrange trajectory, reducing cost and manufacturing error. The limitations of LEO would be the highest traffic level and higher atmospheric drag, which perturbates the satellite slowly into de-orbiting  (useful at the end of life-cycle). However, traffic level in low earth orbit has been an increasingly greater constraint that the Hubble did not have - The low earth orbit region has become increasingly crowded, with more than 8,000 satellites currently orbiting in LEO \cite{UntetherAI2024}, in addition to 29,000 large debris \cite{ESASpaceDebris}. The traffic in front of the telescope may also introduce disturbances to the final results of the image. These are trades-off consequence as a result of lower economical cost and higher servicing accessibility, cons that the original Hubble faces [N6]. 

\subsubsection{Camera}
To enhance the imaging capabilities of ORACLE, the telescope will incorporate several advanced improvements to its optical components, particularly the mirror and sensor systems. The mirror diameter of ORACLE will be further expanded, thereby increasing its light-collecting area and enhancing resolution and clarity for distant celestial objects. This larger mirror will provide more precise details, effectively reducing imaging noise and improving the signal-to-noise ratio, especially when observing galaxies, planets, and stars [N1]. To optimize surface area and manage physical constraints during launch and deployment, the mirror may utilize segmentation techniques similar to those of the James Webb Space Telescope. Furthermore, ORACLE will integrate advanced multi-band sensors, covering a wide spectrum from infrared (IR) to visible light (V) and ultraviolet (UV), allowing simultaneous capture of data across different bands, thus supporting more comprehensive observation tasks [N2,N8,N9]. This enhancement enables the telescope to capture dynamic processes, such as star formation and interstellar gas movements, with greater contrast and detail [N10]. Additionally, to minimize distortion caused by satellite movement and environmental conditions in Low Earth Orbit (LEO), the telescope will be equipped with adaptive optics and active stabilization systems, which will adjust in real-time to maintain image quality, ensuring sharper images even under slight shifts or vibrations [N1].
\vspace{1.5mm}

The field of view (FoV) is set as a requirement [S5] and it serves as the basis the consequent calculations. FoV of 2.7 arcminutes is equivalent to 
\begin{equation}
\text{FoV}_{\text{radians}} = 2.7 \times \frac{1}{60} \times \frac{\pi}{180} \approx 7.85 \times 10^{-4} \text{ radians}
\end{equation}
To achieve this field of view, we calculate the focal length \textit{f} assuming a sensor width that matches this FoV. If the sensor width is 36.2 mm, the focal length required can be estimated to be
\begin{equation}
f \approx \frac{\text{sensor width}}{\text{FoV}} = \frac{36.2 \, \text{mm}}{7.85 \times 10^{-4}} = 46114.6\, \text{mm} = 46.11 \, \text{m}
\end{equation}

We start by selecting the optical system's appropriate wavelength for imaging applications to determine the optimal aperture diameter for achieving the specified angular resolution. This wavelength directly affects the system’s resolution capabilities. The required angular resolution $\theta$, specified in arcseconds, is then converted to radians 
\begin{equation}
\theta_{\text{rad}} = \frac{\theta}{206265}
\end{equation}
And with wavelength and angular resolution defined, the minimum aperture diameter, in m, is calculated 
\begin{equation}
D = 1.22 \cdot \frac{\lambda}{\theta_{\text{rad}}} \ \Rightarrow \ 2.5\leq D\leq 7 
\end{equation}

\vspace{1.5mm}

ORACLE satellite will utilize novel imaging technology, surpassing its predecessor with enhanced capabilities inspired by advancements seen in the James Webb Space Telescope (JWST). This upgrade aims to significantly improve resolution, sensitivity, and overall performance across a wide spectral range, from ultraviolet (UV) to mid-infrared (MIR). A key upgrade will be the use of modern CMOS sensors, which are more efficient and durable compared to the original Hubble’s CCDs, especially in harsh radiation environments. These sensors will cover an extended range, capturing data from infrared (IR) to UV, allowing observations of dust-obscured regions, distant galaxies, and energetic events like supernovas.

\vspace{1.5mm}

 Additionally, ORACLE will also integrate adaptive optics to correct distortions caused by satellite motion and thermal fluctuations. Using techniques similar to those in JWST, real-time adjustments of segmented mirrors will ensure consistently sharp images. These adaptive systems improve the precision of observations, especially for deep-space imaging where clarity is paramount. Moreover, To further enhance sensitivity, the telescope will employ superconducting nanowire single-photon detectors (SNSPDs). These detectors are capable of high quantum efficiency, detecting faint light from distant celestial objects with minimal noise, thus pushing the limits of what can be observed.

 \vspace{1.5mm}

Infrared imaging will be optimized through cryogenic cooling systems, reducing thermal noise and improving signal clarity for exoplanet studies and dust clouds \cite{gardner2006james}. These cooling systems, paired with efficient thermal insulation, will enhance the lifespan and performance of sensitive instruments in the challenging space environment. In terms of data handling, ORACLE will feature high-speed laser communication systems, inspired by recent advancements in optical data transmission, providing faster and more secure data links compared to traditional radio communication. This will enable rapid, high-bandwidth data transfer to Earth, essential for managing the vast amount of data generated by modern sensors.

\vspace{1.5mm}

By integrating these advanced technologies, ORACLE will align with next-generation observatories, significantly boosting its ability to explore distant cosmic phenomena, from the early universe to the study of exoplanetary systems.


\vspace{1.5mm}

\subsubsection{Control \& Communications Features}
ORACLE shall utilize a hemispherical resonator gyroscope (HRG) and a Fine Guidance Sensor instead of a traditional mechanically spinning gyroscope. This offers greater reliability and performance as the current Hubble halted operation due to faulty gyroscope in November 2024\cite{NASANEWS}. The HRG \& FGS is designed to be serviced for at least 15+ year \cite{NASAJWSTFS} and can measure greater precision within Sub-microarcsecond precision for ultra-stable pointing \cite{NASAJWST}. With addition to Secondary Combustion Augmented Thrusters (SCAT), both trajectory and orientational control can be established, allowing the user to take desired/request images with it's corresponding coordinates. [N3][N7] The satellite shall establish telecommunication with a ground station on earth, which allows the user to interact with it remotely using developed software. Command and control can also be transmitted from earth. [N5][N1]

\subsubsection{Sustainability}
Environmental and sustainability of the spacecraft can be achieved with ease, where the end of life (de-orbit) will enter the atmosphere and perform burn out, disintegrating the satellite. Sustainable material and practices shall be practised (More on sustainability on concept 3). [N4]




